\section{Conclusion}
\label{sec:conclusion}
We characterized a structural source of privacy leakage in multi-agent LLM systems: handoff compression preserves operational facts while stripping the boundary metadata that licenses their use. On both GPT-5-mini and DeepSeek-R1-32B, $\sigma$ and $\sigma_{\mathrm{op}}$ decouple at the individual handoff level (Pearson $r$ near zero), and the gap widens under aggressive 25-word compression. Matched-variant downstream tests show that the load-bearing axis is not marker presence but \emph{explicitness}: vague hedges like ``use discretion'' remain high-leakage ($73\%$ on GPT-5-mini, $50\%$ on DeepSeek), while explicit constraints reduce leakage to under $15\%$ across all three tested models. A no-handoff single-agent control further shows that the failure is not reducible to multi-agent topology alone.

These findings reframe what it would take to make multi-agent coordination safe: not better prose summaries, not politer privacy instructions, and not exact-string redaction, but \emph{operationalized boundaries} that downstream agents act on as constraints rather than read as suggestions. A gold-derived typed audience-allowlist schema reduces leakage to $0/48$ on both primary models and far below the no-marker control on Qwen3-32B; a corrupted-allowlist stress test confirms that the load-bearing factor is the correctness and machine-readability of the boundary field, not schema use alone. We release the testbed, evaluator, and analysis scripts to support that work.

Three concrete next steps follow: \emph{automated typed-allowlist extraction} from upstream interactions, with paraphrase-robust validators for category-level leaks redaction misses; \emph{production-trace validation and user studies} of deployed handoffs; and \emph{multilingual replication} of the E8 dose-response on languages with different modal-verb systems.